\section{Protocol}
\label{sec:protocol}

\subsection{Factorial design on the slice}
On the DIPSER slice we cross fusion $\in\{$gate, core$\}$, link $\in\{$on, off$\}$, dropout $\in\{$on, off$\}$ and objective $\in\{$CORAL, hybrid, soft$\}$. This gives 24 fusion configurations, to which we add eight single-stream baselines (four streams, CORAL and soft objectives) and one permuted-label control. With five folds and three seeds this amounts to 495 fits. The control is core fusion with link and dropout on and the CORAL objective, trained on labels permuted across the population; it is required to satisfy $|\mathrm{QWK}|\le0.10$. DAiSEE uses the same 24 configurations and three single-stream baselines (16 distinct fusion models, since soft and CORAL are identical). Because of the time limit of the computing platform, EngageNet uses a reduced set of configurations: three single streams, gate and core fusion with link on and dropout off, and the control (90 fits).

\subsection{Splits}
The main protocol (P2) is a five-fold cross-validation in which no grouping unit appears in both the training and the test fold; QWK is computed on the pooled out-of-fold predictions. On the slice the grouping unit is the student-session. The three students who have complete windows in both sessions may therefore contribute one session to training and the other to testing (5--20\% of the scored windows, depending on the seed); Section~\ref{sec:threats} reports a re-evaluation on person-disjoint windows. On the full release, folds group students, and a second protocol holds out one of the three class groups in turn (leave-one-class-out). For DAiSEE, protocol P1 trains on the official training subjects that do not appear in Validation and tests on Validation (1{,}429 clips); P1-Test retrains the best P1 single-stream and fusion models, selected on Validation, on Train and Validation, and scores them once on Test.

\subsection{Training}
Models are trained with Adam~\cite{kingma2015adam} (learning rate $3\times10^{-3}$, batch size 32, gradient clipping at 1.0) for at most 40 epochs, with early stopping after eight epochs without improvement and not before epoch 9. Epoch selection is nested: 20\% of the subjects of the training fold form an inner validation set, and the checkpoint with the best inner QWK is applied once to the outer fold, which avoids the selection bias of choosing the epoch on the scored fold~\cite{cawley2010overfitting}. A seed fixes the folds, the inner split, the initialisation and the sample order.

\subsection{Statistics}
On the slice, a factorial contrast compares matched configurations that differ in a single factor, with a two-sided Wilcoxon signed-rank test~\cite{wilcoxon1945} on configuration-level means (12 pairs for link, gate and dropout; 8 for the objectives); per-seed tests are reported as secondary. Values written $a\pm b$ are the mean and the standard deviation over three seeds. On the full release, headline values carry 95\% confidence intervals from a bootstrap over students, differences are paired on the same students, and families of related tests are corrected with the Holm procedure. The majority predictor (QWK $=0$) is the reference throughout.

\subsection{Predefined criteria and record of the study}
For the full release, the success criterion of each experiment was written in a dated analysis plan before the corresponding run; the plan, with its amendments, is part of the released material. One experiment, the hyper-parameter search of Section~\ref{sec:fc-limits}, had no predefined criterion and is labelled exploratory. Every reported number is produced by a notebook on a public computing platform (Kaggle), and a separate notebook recomputes the statistics from the stored outputs.
